\documentclass[11pt]{article}

\usepackage[margin=1in]{geometry}
\usepackage{amsmath,amssymb,booktabs,tabularx}

\title{Affine transform ansatz}
\author{}
\date{}

\begin{document}
\maketitle

\noindent\textbf{Status.} Draft\\
\textbf{Plan item.} Affine transform ansatz\\
\textbf{Depends on.} Conditional trading block, news transform, long-variance normalization\\
\textbf{Feeds.} Riccati system, affine domain constraints, terminal pricing measure

\section{Purpose}

State the exponential-affine form used for interval transforms when market
pressure is a finite-state signed-balance CTMC.

\section{Setup}

Let \(\mathbf Y_t\) be the Hawkes reaction state, \(Z_t\) the long-variance
state, and \(\Omega_t\in\{1,\ldots,K_\Omega\}\) the pressure state.
The pressure support value in state \(k\) is \(b_k\in[-1,1]\), the physical
pressure generator is \(A^\Omega\), and the non-negative signed-balance
exposure scale on the interval is \(e^\Omega\).

For a transform return loading \(r_\Omega\) and pressure exposure risk price
\(\rho_\Omega\), define the pressure diagonal loading
\[
    \ell_k^\Omega
    =
    (r_\Omega\gamma_\Omega-\rho_\Omega)e^\Omega b_k.
\]
For physical moment and simulation transforms, \(\rho_\Omega=0\).

\section{Ansatz}

For terminal pressure values \(\mathbf f^\Omega(0)\), the interval transform
has the form
\[
    F_k(\tau,\mathbf y,z)
    =
    \exp\!\left(C(\tau)+\mathbf D(\tau)\cdot\mathbf y+V(\tau)z\right)
    f_k^\Omega(\tau),
\]
where \(k\) is the pressure state at the interval start in backward time.
For an initial pressure distribution \(\mathbf p^\Omega_0\), the pressure
factor is \(\mathbf p^\Omega_0\cdot\mathbf f^\Omega(\Delta)\).

\section{Pressure factor}

The finite-state pressure factor evolves by the backward Kolmogorov equation
\[
    \partial_\tau\mathbf f^\Omega(\tau)
    =
    \left(A^\Omega+\operatorname{diag}(\ell_1^\Omega,\ldots,\ell_{K_\Omega}^\Omega)\right)
    \mathbf f^\Omega(\tau).
\]
Thus
\[
    \mathbf f^\Omega(\Delta)
    =
    \exp\!\left(
        \Delta\left(A^\Omega+\operatorname{diag}(\boldsymbol\ell^\Omega)\right)
    \right)\mathbf f^\Omega(0).
\]

\section{Separation}

The Hawkes and long-variance affine terms remain exponential-affine because the
pressure state changes side composition and anchor drift but does not add a
Hawkes jump loading. The pressure contribution is finite-dimensional and
multiplies the usual affine factor rather than appearing as a scalar loading in
the Hawkes event Hamiltonians.

\section{Checks}

\begin{tabularx}{\textwidth}{XXX}
\toprule
Check & Expected result & Evidence location\\
\midrule
Finite-state pressure factor
& Pressure transform is a matrix exponential over support states.
& \texttt{market\_pressure\_transform\_values}\\
Initial pressure distribution
& Boundary probabilities weight pressure transform values.
& \texttt{market\_pressure\_transform\_log\_factor}\\
Pressure-risk tangent
& The derivative is the diagonal exposure loading \(-e^\Omega b_k\).
& \texttt{market\_pressure\_transform\_values\_and\_derivative}\\
\bottomrule
\end{tabularx}

\end{document}
